% 原创教学示意；三种任务共享节点表示，各自定义读出。
\begin{tikzpicture}[x={\dimexpr\linewidth/169\relax},y=1mm,
  font=\figurefont\footnotesize,text=NNInk]
  \node[nn hidden] (h1) at (12,16) {$\bm h_1$};
  \node[nn hidden] (h2) at (12,0) {$\bm h_2$};
  \node[nn hidden] (h3) at (12,-16) {$\bm h_3$};
  \draw[nn link] (h1)--(h2)--(h3);
  \node[nn operator,minimum width=28mm] (node) at (83,23) {节点分类头};
  \node[nn operator,minimum width=28mm] (pair) at (83,0) {节点对评分};
  \node[nn operator,minimum width=22mm] (pool) at (48,-30) {对称聚合};
  \node[nn operator,minimum width=28mm] (graph) at (100,-30) {整图预测头};
  \draw[nn operator path] (h1.east)--(35,16)--(35,23)--(node.west);
  \draw[nn operator path] (h1.east)--(35,16)--(pair.west);
  \draw[nn operator path] (h2.east)--(pair.west);
  \draw[nn operator path] (h1.west)--(0,16)--(0,-30)--(pool.west);
  \draw[nn operator path] (h2.west)--(0,0);
  \draw[nn operator path] (h3.west)--(0,-16);
  \foreach \yy in {0,-16}{\fill[NNHidden] (0,\yy) circle[radius=.45mm];}
  \draw[nn operator path] (pool)--node[above] {$\bm g$} (graph);
  \node[anchor=west,text=NNOutput] (pn) at (126,23) {$\bm p_1$};
  \node[anchor=west,text=NNOutput] (pe) at (126,0) {$p_{1,2}$};
  \node[anchor=west,text=NNOutput] (pg) at (126,-30) {$\widehat y$};
  \draw[nn flow,draw=NNOutput] (node)--(pn);
  \draw[nn flow,draw=NNOutput] (pair)--(pe);
  \draw[nn flow,draw=NNOutput] (graph)--(pg);
\end{tikzpicture}
